\paragrafos{Δεκαδική μορφή ρητών: όταν τα ψηφία δεν τελειώνουν}{Α.7.7, «Δεκαδική μορφή ρητών αριθμών»}

\begin{anagke}
Ο Στέλιος παίζει με το κομπιουτεράκι του κινητού και πατάει $1:3$. Βλέπει $0{,}33333333$ και συνεχίζει με το $1:7$: $0{,}14285714$. «Το $1:7$ έχει ένα σωρό ψηφία!». Η Πολυξένη το αντιγράφει και προσπαθεί να δει αν υπάρχει κάποιο μοτίβο: «$142857\ 14\ldots$ Το $142857$ επαναλαμβάνεται!». Δοκιμάζει με το $2:7$ και βρίσκει $0{,}285714\ 285\ldots$: ίδια ψηφία, αλλά ξεκινούν από άλλο σημείο.

Ο Νίκος δοκιμάζει $3:8$ και βρίσκει $0{,}375$ και τελειώνει. Η Αγγελική ρωτάει: «Πότε τελειώνει η διαίρεση και πότε όχι;». Η Μυρτώ κάνει την ευκλείδεια διαίρεση στο χέρι: «Στο $1:7$ τα υπόλοιπα είναι πάντα $1,2,\dots,6$, δεν έχουν άλλες επιλογές. Μετά από έξι βήματα θα ξαναβρούμε κάποιο υπόλοιπο και τα ψηφία θα επαναληφθούν».

Και ο Βασίλης προτείνει το δύσκολο: «Αν πάρω το $0{,}3333\ldots$ και το πολλαπλασιάσω επί τρία, βγαίνει $0{,}9999\ldots$. Αλλά το $\frac{1}{3}\cdot3=1$. Άρα το $0{,}9999\ldots$ ισούται με το $1$;». Σε αυτή την παράγραφο η παρέα βρίσκει τον κανόνα του πότε τελειώνει ένα δεκαδικό ανάπτυγμα, μαθαίνει να γράφει τους άπειρους δεκαδικούς και πώς να γυρίζει πίσω στο κλάσμα.
\end{anagke}

\begin{lexi}{περιοδικός, περίοδος}{περιοδικός, περίοδος}
\lexiword{Περίοδος} από το «περί-οδος», δηλαδή «γύρω-δρόμος»: ένας δρόμος που τον ξανακάνουμε. Ο \lexiword{περιοδικός δεκαδικός} έχει ψηφία που επαναλαμβάνονται σε κύκλο (την περίοδο) για πάντα. Ο \lexiword{ακριβής δεκαδικός} τελειώνει, χωρίς να επαναλαμβάνεται τίποτα.
\end{lexi}

\begin{metafrasi}
Το $\frac{3}{8}$ τελειώνει. & $0{,}375$\\
Το $\frac{1}{3}$ επαναλαμβάνει το $3$. & $0{,}\overline{3}=0{,}333\ldots$\\
Το $\frac{1}{7}$ επαναλαμβάνει το $142857$. & $0{,}\overline{142857}$\\
Το $\frac{5}{6}$: το $8$ μία φορά και μετά το $3$. & $0{,}8\overline{3}=0{,}8333\ldots$\\
Από περιοδικό σε κλάσμα. & $0{,}\overline{27}=\dfrac{27}{99}=\dfrac{3}{11}$\\
\end{metafrasi}

\tmima{Η θεωρία}

Κάθε ρητός $\frac{\kappa}{\lambda}$ γίνεται δεκαδικός αν κάνουμε τη διαίρεση $\kappa:\lambda$. Υπάρχουν δύο ενδεχόμενα:

\begin{itemize}[leftmargin=1.2em,itemsep=0.3ex,topsep=2pt]
\item Η διαίρεση \textbf{τελειώνει} (υπόλοιπο $0$): έχουμε \textbf{ακριβή δεκαδικό} (πεπερασμένο πλήθος δεκαδικών ψηφίων). Π.χ. $\frac{3}{8}=0{,}375$, $\frac{9}{20}=0{,}45$.
\item Η διαίρεση \textbf{δεν τελειώνει}: τα υπόλοιπα είναι πάντα μικρότερα από τον διαιρέτη $\lambda$, άρα κάποτε επαναλαμβάνονται και μαζί τους επαναλαμβάνονται τα ψηφία του πηλίκου. Έχουμε \textbf{περιοδικό δεκαδικό}. Γράφουμε $\frac{5}{6}=0{,}8\overline{3}$ και $\frac{4}{11}=0{,}\overline{36}$.
\end{itemize}

Το τμήμα που επαναλαμβάνεται λέγεται \textbf{περίοδος}. Αν η περίοδος αρχίζει αμέσως μετά την υποδιαστολή, ο δεκαδικός λέγεται \emph{απλός περιοδικός} ($0{,}\overline{36}$), αλλιώς \emph{μικτός} ($0{,}8\overline{3}$).

\begin{kanonas}{πότε τελειώνει}
Ένα ανάγωγο κλάσμα γράφεται ως ακριβής δεκαδικός όταν ο παρονομαστής του έχει \textbf{μόνο παράγοντες $2$ και $5$} (γιατί τότε γίνεται δεκαδικό κλάσμα με παρονομαστή δύναμη του $10$). Αλλιώς είναι περιοδικός.
\end{kanonas}

Έτσι το $\frac{7}{20}$ ($20=2^2\cdot5$) τελειώνει: $0{,}35$. Το $\frac{5}{12}$ ($12=2^2\cdot3$) δεν τελειώνει: $0{,}41\overline{6}$.

\tmimasmall{Από περιοδικό σε κλάσμα}

Αν $x=0{,}\overline{3}$, τότε $10x=3{,}\overline{3}$ και αφαιρώντας $10x-x=3$, άρα $9x=3$ και $x=\frac{1}{3}$. Με την ίδια ιδέα:
\begin{itemize}[leftmargin=1.2em,itemsep=0.3ex,topsep=2pt]
\item \textbf{Απλός περιοδικός:} αριθμητής η περίοδος και παρονομαστής τόσα εννιάρια όσα τα ψηφία της περιόδου: $0{,}\overline{27}=\frac{27}{99}=\frac{3}{11}$, $0{,}\overline{7}=\frac{7}{9}$, $1{,}\overline{3}=1+\frac{3}{9}=\frac{4}{3}$.
\item \textbf{Μικτός:} $0{,}2\overline{3}$: με $x=0{,}2\overline{3}$, $10x=2{,}\overline{3}=2+\frac{1}{3}=\frac{7}{3}$, άρα $x=\frac{7}{30}$.
\end{itemize}
Κάθε περιοδικός δεκαδικός είναι ρητός, και κάθε ρητός είναι ακριβής ή περιοδικός δεκαδικός. Ο $0{,}\overline{9}=1$: με $x=0{,}\overline{9}$, $10x=9{,}\overline{9}$, $9x=9$, $x=1$.

\begin{prosoxi}
\begin{itemize}[leftmargin=1.2em,itemsep=0.3ex,topsep=2pt]
\item Το $\frac{1}{3}$ δεν είναι ίσο με $0{,}33$ αλλά με $0{,}\overline{3}$: το $0{,}33$ είναι προσέγγιση.
\item Η περίοδος είναι μόνο το κομμάτι που επαναλαμβάνεται: στο $0{,}8\overline{3}$ το $8$ δεν επαναλαμβάνεται.
\item Η περίοδος ενός κλάσματος με παρονομαστή $\lambda$ έχει το πολύ $\lambda-1$ ψηφία (λόγω των υπολοίπων).
\item Τα κομπιουτεράκια κόβουν ή στρογγυλοποιούν τα ψηφία: $1:7$ δεν σταματά στο $0{,}14285714$.
\end{itemize}
\end{prosoxi}

\tmima{Λυμένα παραδείγματα}

\paradeigma{Ακριβής ή περιοδικός;}
\tixero Τα κλάσματα $\frac{7}{20}$, $\frac{5}{12}$, $\frac{9}{14}$, $\frac{11}{40}$.\par
\tithelo Ποια έχουν ακριβή δεκαδικό.\par
\possindeo Αναλύω τους παρονομαστές: $20=2^2\cdot5$, $12=2^2\cdot3$, $14=2\cdot7$, $40=2^3\cdot5$. Μόνο παράγοντες $2$ και $5$ έχουν το $20$ και το $40$.\par
\apantisi Ακριβείς: $\frac{7}{20}=0{,}35$ και $\frac{11}{40}=0{,}275$. Περιοδικοί: $\frac{5}{12}=0{,}41\overline{6}$ και $\frac{9}{14}=0{,}6\overline{428571}$.\par
\elegxos Το $0{,}35=\frac{35}{100}=\frac{7}{20}$.

\paradeigma{Η διαίρεση του 1:7}
\tixero Το $\frac{1}{7}$.\par
\tithelo Να βρω τα ψηφία και την περίοδο.\par
\possindeo $1:7$: $10:7=1$ υπόλοιπο $3$, $30:7=4$ υπόλοιπο $2$, $20:7=2$ υπόλοιπο $6$, $60:7=8$ υπόλοιπο $4$, $40:7=5$ υπόλοιπο $5$, $50:7=7$ υπόλοιπο $1$. Το υπόλοιπο $1$ ξαναεμφανίζεται, άρα η περίοδος είναι $142857$.\par
\apantisi $\frac{1}{7}=0{,}\overline{142857}$.\par
\elegxos Η περίοδος έχει $6$ ψηφία, όσα τα δυνατά υπόλοιπα $1,\dots,6$.

\paradeigma{Από περιοδικό σε κλάσμα}
\tixero Οι $0{,}\overline{27}$, $2{,}\overline{6}$ και $0{,}1\overline{6}$.\par
\tithelo Να τους γράψω ως κλάσματα.\par
\possindeo $0{,}\overline{27}=\frac{27}{99}=\frac{3}{11}$. $2{,}\overline{6}=2+\frac{6}{9}=2+\frac{2}{3}=\frac{8}{3}$. Για τον τρίτο $x=0{,}1\overline{6}$, $10x=1{,}\overline{6}=1+\frac{2}{3}=\frac{5}{3}$, άρα $x=\frac{5}{30}=\frac{1}{6}$.\par
\apantisi $\frac{3}{11}$, $\frac{8}{3}$, $\frac{1}{6}$.\par
\elegxos Το $\frac{1}{6}=0{,}1666\ldots$ και το $3:11=0{,}2727\ldots$.

\paradeigma{Στρογγυλοποίηση περιοδικών}
\tixero Οι $\frac{2}{3}$, $\frac{5}{6}$, $\frac{1}{7}$.\par
\tithelo Στρογγυλοποίησε στα εκατοστά.\par
\possindeo $\frac{2}{3}=0{,}666\ldots\approx0{,}67$. $\frac{5}{6}=0{,}8333\ldots\approx0{,}83$. $\frac{1}{7}=0{,}1428\ldots\approx0{,}14$.\par
\apantisi $0{,}67$, $0{,}83$, $0{,}14$.\par
\elegxos Το σφάλμα στο πρώτο: $0{,}67-0{,}666\ldots=0{,}0033\ldots<0{,}005$.

\paradeigma{Ισούται το $0{,}\overline{9}$ με το $1$;}
\tixero Ο αριθμός $0{,}\overline{9}$.\par
\tithelo Να τον συγκρίνω με το $1$.\par
\possindeo Αν $x=0{,}\overline{9}$, τότε $10x=9{,}\overline{9}$. Αφαιρώ: $10x-x=9$, άρα $9x=9$ και $x=1$. Ή: $\frac{1}{3}=0{,}\overline{3}$, άρα $3\cdot\frac{1}{3}=3\cdot0{,}\overline{3}=0{,}\overline{9}$, δηλαδή $1=0{,}\overline{9}$.\par
\apantisi $0{,}\overline{9}=1$.\par
\elegxos Δεν υπάρχει αριθμός ανάμεσα στο $0{,}\overline{9}$ και στο $1$.

\tmima{Ασκήσεις}

\zoni{Το κατάλαβα}{ακριβείς και περιοδικοί}

\ex Γράψε σε δεκαδική μορφή (ακριβή ή περιοδική): α) $\frac{3}{8}$ \quad β) $\frac{5}{6}$ \quad γ) $\frac{4}{11}$ \quad δ) $\frac{7}{9}$ \quad ε) $\frac{9}{20}$ \quad στ) $\frac{2}{15}$ \quad ζ) $\frac{1}{12}$.

\ex Ποια από τα $\frac{7}{20}$, $\frac{5}{12}$, $\frac{9}{14}$, $\frac{11}{40}$, $\frac{3}{16}$, $\frac{13}{25}$ γράφονται ως ακριβείς δεκαδικοί; Χωρίς να κάνεις τη διαίρεση.

\ex Γράψε τα πρώτα οκτώ δεκαδικά ψηφία των $0{,}\overline{36}$, $2{,}1\overline{45}$ και $0{,}\overline{142857}$.

\ex Γράψε ως κλάσματα: α) $0{,}\overline{7}$ \quad β) $0{,}\overline{15}$ \quad γ) $1{,}\overline{3}$ \quad δ) $0{,}\overline{81}$ \quad ε) $0{,}2\overline{3}$.

\ex Σύγκρινε: α) $0{,}\overline{3}$ και $0{,}33$ \quad β) $0{,}\overline{6}$ και $\frac{2}{3}$ \quad γ) $0{,}7\overline{2}$ και $0{,}72$.

\ex Στρογγυλοποίησε στα εκατοστά: α) $\frac{2}{3}$ \quad β) $\frac{5}{6}$ \quad γ) $\frac{1}{7}$ \quad δ) $\frac{4}{9}$.

\ex Σωστό ή λάθος; α) Κάθε ρητός είναι ακριβής ή περιοδικός δεκαδικός. \quad β) $\frac{1}{3}=0{,}33$. \quad γ) $0{,}\overline{9}=1$. \quad δ) Το $\frac{1}{7}$ έχει περίοδο $142857$. \quad ε) Τα κλάσματα με παρονομαστή $8$ είναι ακριβείς δεκαδικοί. \quad στ) $\frac{1}{6}=0{,}16$.

\zoni{Το χρησιμοποιώ}{κομπιουτεράκι, μοιρασιές, μέσοι όροι}

\ex \textbf{Το κομπιουτεράκι.} Ο Στέλιος γράφει $1:7$ και βλέπει $0{,}14285714$. Ποια είναι η περίοδος; Βρες με το κομπιουτεράκι το $2:7$ και το $3:7$ και πες τι παρατηρείς.

\ex \textbf{Η μοιρασιά.} Τα $10$ € μοιράζονται σε $3$ παιδιά. Τι δείχνει το $10:3$ και πόσο παίρνει στην πράξη το καθένα;

\ex \textbf{Η ταχύτητα.} Ο Νίκος τρέχει $100$ μ. σε $12$ s. Πόση είναι η ταχύτητά του σε m/s ως περιοδικός και ως στρογγυλοποιημένος στα εκατοστά;

\ex \textbf{Τα ποσοστά.} Γράψε το $\frac{1}{3}$ και το $\frac{5}{6}$ ως ποσοστά με περιοδικούς και στρογγυλοποίησε στο δέκατο.

\ex \textbf{Ο μέσος όρος.} Οι βαθμοί $7$, $8$, $8$ έχουν μέσο όρο $\frac{23}{3}$. Γράψε τον ως περιοδικό και στρογγυλοποίησέ τον στα εκατοστά.

\zoni{Το σκέφτομαι}{γιατί επαναλαμβάνονται}

\ex Εξήγησε γιατί η διαίρεση $\kappa:7$ έχει περίοδο το πολύ $6$ ψηφίων.

\ex Δείξε ότι $0{,}\overline{9}=1$ με τον τρόπο $x=0{,}\overline{9}$, $10x=9{,}\overline{9}$.

\ex Βρες τη δεκαδική μορφή του $\frac{1}{11}$ και του $\frac{1}{13}$ (δοκίμασε με κομπιουτεράκι και με διαίρεση).

\ex Βρες έναν ρητό ανάμεσα στους $0{,}\overline{3}$ και $0{,}34$.

\ex Γιατί το $\frac{1}{3}$ ισούται με $0{,}\overline{3}$ και όχι με $0{,}333$; Τι μας λέει το $3\cdot0{,}\overline{3}$;

\zoni{Για το σπίτι}{ανακάλυψε τα μοτίβα}

\ex Υπολόγισε με κομπιουτεράκι το $1:\nu$ για $\nu=2,\dots,12$ και βρες ποια έχουν ακριβές ανάπτυγμα. Τι κοινό έχουν οι παρονομαστές τους;
